\begin{theorem}[逆估计]\label{thm:appendix-c-30}
设 $0\le k\le l$ 为整数,且 $p,q\in[1,\infty]$.则存在常数 $C_{\mathrm{inv}}$,它只依赖于 $k,l,p,q,\widehat K$ 及 $\widehat P(\widehat K)$,使得
\begin{equation}\label{eq:appendix-c-35}
 \|D^l v^h\|_{L^q(K)}
 \le C_{\mathrm{inv}}h_K^{(k-l)-d(p^{-1}-q^{-1})}
       \|D^k v^h\|_{L^p(K)},
 \qquad \forall v^h\in P(K).
\end{equation}
\end{theorem}

\begin{proof}
首先在参考网格单元上证明式\hyperref[eq:appendix-c-35]{(C.35)},其中
$h_{\widehat K}=1$ 且 $k=0$.由于有限维空间中任意两种范数等价,存在常数 $C$,使得
\begin{equation}\label{eq:appendix-c-36}
 \|D^l\widehat v^h\|_{L^q(\widehat K)}
 \le\|\widehat v^h\|_{W^{l,q}(\widehat K)}
 \le C\|\widehat v^h\|_{L^p(\widehat K)},
 \qquad \forall\widehat v^h\in\widehat P(\widehat K).
\end{equation}
若 $k>0$,则令
\[
 \widetilde P(\widehat K)
 =\{\partial^\alpha\widehat v^h:
    \widehat v^h\in\widehat P(\widehat K),\ |\alpha|=k\}.
\]
这仍是由多项式组成的空间.将式\hyperref[eq:appendix-c-36]{(C.36)} 应用于
$\widetilde P(\widehat K)$,得到
\[
\begin{aligned}
 \|D^l\widehat v^h\|_{L^q(\widehat K)}
 &=\sum_{|\alpha|=k}
    \|D^{l-k}\partial^\alpha\widehat v^h\|_{L^q(\widehat K)}\\
 &\le C\sum_{|\alpha|=k}
    \|\partial^\alpha\widehat v^h\|_{L^p(\widehat K)}
 =C\|D^k\widehat v^h\|_{L^p(\widehat K)}.
\end{aligned}
\]
按定理\hyperref[thm:appendix-c-13]{C.13} 证明中的相同方法,将该估计变换到一般网格单元 $K$ 上.由变换估计可得
\[
\begin{aligned}
 \|D^l v^h\|_{L^q(K)}
 &\le Ch_K^{-l+d/q}
       \|D^l_{\widehat{\bm x}}\widehat v^h\|_{L^q(\widehat K)}\\
 &\le Ch_K^{-l+d/q}
       \|D^k_{\widehat{\bm x}}\widehat v^h\|_{L^p(\widehat K)}\\
 &\le C_{\mathrm{inv}}h_K^{k-l+d/q-d/p}
       \|D^k v^h\|_{L^p(K)}.
\end{aligned}
\]
注意 $k-l+d/q-d/p=(k-l)-d(p^{-1}-q^{-1})$,故得式\hyperref[eq:appendix-c-35]{(C.35)}.
\end{proof}

\begin{note}[关于证明]\label{note:appendix-c-31}
证明中的关键是有限维空间中所有范数彼此等价.无限维空间并不具有这一性质.
\end{note}

\begin{corollary}[全局逆估计]\label{cor:appendix-c-32}
设 $p=q$,且 $\mathcal T^h$ 是 $\Omega$ 上的准一致剖分族,则
\begin{equation}\label{eq:appendix-c-37}
 \|D^l v\|_{L^p_h(\Omega)}
 \le C_{\mathrm{inv}}h^{k-l}
       \|D^k v\|_{L^p_h(\Omega)}.
\end{equation}
其中
\[
 \|v\|_{L^p_h(\Omega)}
 =\left(\sum_{K\in\mathcal T^h}\|v\|_{L^p(K)}^p\right)^{1/p}.
\]
\end{corollary}

\begin{note}[关于 $\|\cdot\|_{L^p_h(\Omega)}$]\label{note:appendix-c-33}
按网格单元定义该范数十分重要:当 $k\ge2$ 或 $l\ge2$ 时,有限元函数通常不具备全局范数有定义所需的正则性;当 $l\ge1$ 且有限元函数非协调时,这种网格单元范数也很重要.
\end{note}
